\documentclass[11pt]{article}
\usepackage[T1]{fontenc}
\usepackage[utf8]{inputenc}
\usepackage[a4paper,margin=1.72cm]{geometry}
\usepackage{amsmath,amssymb}
\pagestyle{plain}
\setlength{\parskip}{1.0pt}
\newcommand{\A}{\mathcal A}
\newcommand{\G}{\mathcal G}
\newcommand{\Mf}{\mathcal M_{\rm flat}}
\newcommand{\z}{\mathfrak z}

\begin{document}

\begin{center}
{\Large\bfseries Erratum and restoration (v2): one symbol carries two\\
angles, several printed statements are false, (H1) is not\\
discharged, and (H-FACT) does not replace (H2)}\\[5pt]
{\large Morse--Bott Spectral Reduction and the\\
Yang--Mills Mass Gap on the Lattice}\\[6pt]
Lluis Eriksson \quad\textbullet\quad ai.viXra:2602.0035 \quad\textbullet\quad July 2026
\end{center}

\noindent\rule{\textwidth}{0.8pt}

\vspace{2pt}
\noindent
\textbf{This version does two things at once.} It \emph{restores} the revision prepared
for this record in July 2026, which was mis-filed onto \texttt{ai.viXra:2602.0052} and
so never appeared here; and it \emph{corrects} that revision. \textbf{It refutes several
printed statements}, and they do not all fail for the same reason. \textbf{The angle
ambiguity accounts for some, but not all, of the defects}: the symbol $\theta$ carries
both the loop holonomy angle and the per-link background angle, which differ by a factor
$L$ (\S2), and that invalidates equation (5), the shift used in equation (6), the Fourier
labels of Proposition 3.3 and the frequencies (9). \textbf{Proposition 2.2(b), the
additive split in equation (6), the centraliser count and the normal-bundle
classification fail for independent reasons.} Two status claims --- that (H1) is
discharged, and that (H2) is refined into (H-FACT) --- are withdrawn in \S4.

\section{What happened to this record}

The file prepared as v2 of this paper was uploaded, during the replacement campaign of
2026-07-06, onto the record \texttt{ai.viXra:2602.0052}, whose own paper
(\emph{Interface Lemmas}) it displaced. This record therefore still showed v1 while
three companion papers cited ``\texttt{2602.0035} v2'' --- a version that existed only
as the mis-filed file. The pages after this note are that revision, restored; the
corrections below apply to it.

\section{One symbol, two angles}

\S2.1 defines the total holonomy $g_\mu=\prod_{t=0}^{L-1}U_\mu(t\hat\mu)=e^{i\Theta_\mu}$
--- writing $\Theta$ for what the manuscript calls $\theta$. The flat,
translation-invariant representative that the Fourier analysis then uses must have
\[
U_\mu(x)=e^{i\Theta_\mu/L},
\]
because there are $L$ links around the cycle. \textbf{Any corrected proof must
distinguish the loop-holonomy angle from the per-link angle.} Two consequences follow
immediately, and both are printed the other way.

\emph{The covariant symbol.} For a fluctuation in root direction $\alpha$ the eigenvalue
of the covariant difference is
\[
4\sin^2\!\Bigl(\frac{k_\mu+\alpha(\Theta_\mu)/L}{2}\Bigr)
\qquad\text{--- in flat notation, } k + \text{alpha(Theta)/L} \text{ ---}
\]
and not $4\sin^2((k_\mu+\alpha(\Theta_\mu))/2)$.

\emph{The metric.} In total-holonomy coordinates the per-link variation is
$\delta\Theta_\mu/L$, so
$\sum_x|\delta\Theta_\mu/L|^2=L^d\,L^{-2}|\delta\Theta_\mu|^2
=L^{(d-2)}|\delta\Theta_\mu|^2$, whereas equation (5) prints
$L^d|\delta\Theta_\mu|^2$.

Only two readings are available, and each demands a correction. If $\theta$ is the
total-holonomy angle, as Proposition 2.2 says, then the factors $1/L$ are missing from
the metric and from the shifted momentum. If $\theta$ is the per-link angle, it is no
longer the coordinate of Proposition 2.2 and its period is $2\pi/L$, with the
corresponding gauge identifications. \textbf{The manuscript uses one symbol for both
without declaring a change of variable}, so equation (5), the mass term
$m_a^2(\theta)$, the proof of (6) and the frequencies (9) cannot all be correct as
printed.

\section{What is false as printed}

\subsection{Proposition 2.2(b): the $SU(2)$ quotient is misidentified}

The proposition states $\Mf\cong[0,\pi]^d/\mathbb Z_2$ with
$\theta_\mu\mapsto\pi-\theta_\mu$. But the SU(2) Weyl action is simultaneous inversion:
$(\theta_1,\dots,\theta_d)\mapsto(-\theta_1,\dots,-\theta_d)$ on $(S^1)^d$. The printed
proof passes from $\theta\mapsto2\pi-\theta$ to $\theta\mapsto\pi-\theta$ with no change
of coordinates; these are not the same map. The paper contradicts itself on it: under
inversion the fixed points are the $2^d$ tuples with $\theta_\mu\in\{0,\pi\}$, while
under the printed map the only interior fixed point is $(\pi/2,\dots,\pi/2)$ --- and
Proposition 3.3 then places $\theta=0$ on the orbifold locus, which the printed map does
not fix. \textbf{Proposition 2.2(b) is false as printed.} The correct statement is
$\Mf^{SU(2)}=(S^1)^d/\{\theta\sim-\theta\}$; a cube presentation is possible but needs
its boundary identifications fixed with care.

\subsection{Equation (5): the metric factor}

As computed in \S2, the induced metric in total-holonomy coordinates carries $L^{d-2}$,
not $L^d$. \textbf{Equation (5) is false as printed} under the reading Proposition 2.2
fixes.

\subsection{Equation (6): two independent failures}

The theorem asserts
$\operatorname{Hess}(V_{\rm pot})^\perp=\frac{2}{g^2}\sum_{k\neq0,a}
[\widehat k^2+m_a^2(\theta)]\,|\delta A^a_\perp(k)|^2$.

\emph{First failure --- the split.} Its own proof gives the covariant-difference
eigenvalue as $4\sum_\mu\sin^2((k_\mu+\phi_\mu)/2)$, and
\[
4\sin^2\!\Bigl(\frac{k+\phi}{2}\Bigr)\;\neq\;
4\sin^2\!\Bigl(\frac{k}{2}\Bigr)+4\sin^2\!\Bigl(\frac{\phi}{2}\Bigr).
\]
\textbf{A test inside the smooth stratum}: $SU(2)$, $L=6$, $\Theta=\pi/2$, $\alpha=2$,
$k=\pi/3$. Then $\phi=\alpha\Theta/L=\pi/6$, the left side is $4\sin^2(\pi/4)=2$ and
the right side is $3-\sqrt3\approx1.268$. \emph{$\Theta=\pi/2$ is a regular holonomy},
so the test lies where Theorem 3.1 is restricted. \emph{An earlier draft used
$k=\phi=\pi/3$, which under the corrected convention forces $\Theta=\pi$ --- central
holonomy, outside the smooth stratum, and so an avoidable objection about the domain.}
\textbf{Equation (6) is not valid as a general equality away from trivial holonomy}
--- stated that way because special pairs $(k,\phi)$ may agree by accident. The proof
concedes it without saying so: it obtains (6) \emph{at trivial holonomy} and, elsewhere,
replaces the equality by an inequality with a cross term. That inequality may serve an
argument for positivity; it does not establish the displayed formula.

\emph{Second failure, independent of the first --- the shift.} By \S2 the shift is
$\phi_\mu=\alpha(\Theta_\mu)/L$, not $\alpha(\Theta_\mu)$. So even the inequality, and
the mass term $m_a^2$, are written with the wrong argument.

\emph{Why the suite does not catch it.} The Hessian is validated at trivial holonomy,
where the shift vanishes and the defect is therefore absent. This is the same pattern as
the one-dimensional miniature in \texttt{ai.viXra:2602.0033}: \textbf{the test does not
examine the regime in which the defect lives.} Two papers of this chain now have a check
with that property: a fact about how test cases are chosen, not a coincidence.

\subsection{Proposition 3.3: the count, and the Fourier labels}

The proposition counts $d(N_c^2-1-r)$ massless non-Cartan \emph{$k=0$ constant} modes at
any point of enhanced stabiliser. Two things are wrong.

\emph{The count.} Its own mass formula makes a root direction $\alpha$ massless exactly
when its pairing with every holonomy angle is an integer multiple of $2\pi$ --- that is,
exactly on the common centraliser. \textbf{The correct count is $d(\dim\z_\Theta-r)$},
in flat notation $d$ times (dim z\_theta $-$ r), with $\z_\Theta$ the algebra of the
common centraliser; the printed count holds only when that centraliser is the whole
group. For $SU(3)$ at a point with centraliser $S(U(2)\times U(1))$, $\dim\z_\Theta=4$
and $r=2$, giving $2d$, not $6d$.

\emph{The labels.} \textbf{Covariantly constant modes need not have Fourier label k=0.}
Take $SU(2)$ with the central holonomy $g_\mu=-I$, $\Theta_\mu=\pi$, and the adjoint root
$\alpha=2$: the per-link shift is $2\pi/L$, so
$\lambda_\alpha(k=0)=4d\sin^2(\pi/L)>0$, while the zero sits at the allowed momentum
$k_\mu=-2\pi/L$. \textbf{The centraliser is everything here, so this is precisely a case
the proposition claims}, and the modes it claims at $k=0$ are not at $k=0$. The accurate
statement is: at an enhanced-stabiliser holonomy there are $d(\dim\z_\Theta-r)$
additional covariantly constant centraliser directions, and their Fourier labels depend
on the chosen flat representative.

\emph{And ``normal to $\Mf$'' is not a valid classification at a singular point}, where
there is no ordinary normal bundle. At the trivial point the quartic form is
$\propto\sum_{\mu<\nu}|[X_\mu,X_\nu]|^2$, which \emph{vanishes on commuting tuples}:
taking all $X_\mu$ proportional to one non-Cartan generator keeps the connection flat,
with no quartic growth. The accurate classification is that all constant directions lie
in the kernel of the Hessian at the trivial point, commuting tuples lie in the tangent
cone to the flat connections, and non-commuting tuples exhibit the positive quartic term.

What the numerics support is that for $SU(2)$ on a $2^3$ spatial lattice the Hessian
kernel has dimension $30$ and a non-commuting direction has quartic ratio $15.99$.
\textbf{Verified, not proved}, and for that instance.

\subsection{Proposition 4.1: the cutoff does not do what it is asked to do}

Correcting the Fourier labels refutes two further printed claims, and this propagates
into the Born--Oppenheimer reduction rather than stopping at Proposition 3.3.

\emph{The cutoff.} Proposition 4.1 says the sums (8) run over $k\neq0$ and thereby omit
the toron modes. The mode of \S3.4 --- $g_\mu=-I$, $\alpha=2$, $k_\mu=-2\pi/L$ --- has
$k\neq0$, lies \emph{inside} those sums, and has zero frequency. \textbf{The k != 0
cutoff does not exclude all covariantly constant modes.}

\emph{The uniform gap.} The printed normal-mode bound is
$\omega_{\perp,\min}=2\sqrt2\sin(\pi/L)>0$, uniform over the smooth stratum. Put
$\Theta_\mu=\pi-\epsilon$ with $\epsilon\neq0$, so the holonomy is regular, and keep
$k_\mu=-2\pi/L$: then $k_\mu+\alpha\Theta_\mu/L=-2\epsilon/L$, the eigenvalue is
$\lambda_\epsilon=4d\sin^2(\epsilon/L)$ and the corresponding frequency of (9) is
$\omega_\epsilon=\sqrt{2\lambda_\epsilon}=\sqrt{8d}\,|\sin(\epsilon/L)|\to0$ as
$\epsilon\to0$, at \emph{fixed} $L$ --- written as a frequency so that the comparison is
with like for like --- while the printed bound does not move. \textbf{The normal-mode gap closes on approach to the
enhanced-stabiliser locus.} So there is no such uniform bound over the smooth stratum.

\textbf{Both claims of Proposition 4.1 are false as printed.} The printed harmonic
Born--Oppenheimer reduction, \textbf{and any argument relying on a uniformly gapped
normal sector, can at most hold on a region uniformly separated from the
enhanced-stabiliser locus}, with a holonomy-dependent normal gap. \emph{A more general
global theorem is not excluded}: it would additionally require the quartic-sector
construction listed in \S7(C2). Equations (8)--(10), the bound on the remaining
eigenvalues, and the proof's explanatory remark all inherit this.

\subsection{What this leaves of Morse--Bott}

\textbf{The smooth-stratum conclusion is not refuted but is not established} by the
printed proof. Equation (6) is false as an equality, and the shift it uses is wrong by a
factor $L$, so the printed non-degeneracy argument is gone. A corrected proof must
distinguish the two angles and use the actual shifted spectrum; the conclusion is
plausible and may well survive, but this manuscript does not establish it. The
frequencies (9) inherit both defects, in a proposition the text already labels
\emph{formal}.

\section{(H1), (H2), and what Theorem 6.3 is}

\subsection{(H1) is not discharged}

The restored text asserts in four places --- Proposition 6.4, Remark 6.6, item (4) of
the Note added, and \S7.3 --- that (H1) is discharged by \texttt{ai.viXra:2602.0033} v2,
Theorem 3.1. \textbf{All four are withdrawn.} That companion's erratum records that its
Theorem 3.1 is not established by its printed proof, and that its Theorem 4.1 --- the
doubling cascade --- is independently not established, its transfer operator changing
with the isotropic volume over which the limit is taken.

\textbf{Only the negative status assessment of Proposition 6.4 remains accurate}: at the
relevant coupling the terminal gap is open, Balaban's regime being intermediate rather
than strong. Its \emph{positive} appeal to Corollary 5.5 is not available either, since
that corollary rests on Theorem 5.3 against $\mathrm{Vol}_B$ and on the withdrawn Ricci
bound; it would need the corrected physical-measure construction and
$(\text{A-osc})_\nu$ of \S5.4.

Moreover (H1) as printed --- ``\emph{the spectral gap of the transfer matrix associated
with $S_{n_{\max}}$ on the orbit space}'' --- \textbf{does not by itself fix a space, a
measure or an operator}, as the audit of \texttt{2602.0033} showed. A usable (H1) must
name the physical transfer operator, its fixed spatial slice and its Hilbert space.

\subsection{(H-FACT) does not replace (H2)}

Proposition 6.5, Remark 6.6, \S7.3 and Note added (4) present (H2) --- \emph{the
physical gap is RG-invariant} --- as \emph{refined} into (H-FACT), the non-extensive
temporal factorisation hypothesis of the companion. \textbf{That is withdrawn too, and
it is not a matter of wording.} Getting from (H-FACT) to (H2) was the work of Theorem
4.1 and Proposition 4.4 of \texttt{2602.0033}, and that extraction is exactly what
fails. \textbf{H-FACT does not by itself imply H2} with the printed proofs; it is at
most one input to a future proof of (H2). \S7.3, which called (H-FACT) the remaining
analytic content of the programme, is withdrawn with it.

\subsection{Theorem 6.3}

\textbf{The algebraic implication survives on admissible volumes}, under
\emph{Balaban's Theorem 6.1} --- (H-BAL) in the coordinated vocabulary, and the
theorem's own first line assumes it --- together with the original (H1) and (H2), once
those hypotheses are attached to precisely specified operators. \emph{An earlier draft
of this note omitted (H-BAL) from that list.}

\emph{And not with the same prefactor.} \textbf{Equation (18), not the final proof
line, retains the c\_0 of (H1)}: it states $m_{\rm gap}\ge c_0e^{-C/g^2}$ with the very
$c_0$ of that hypothesis, whereas \textbf{the proof already introduces c'\_0}, writing
$m_0\ge c_0e^{-C/g^2-O(1)}=c_0'e^{-C/g^2}$ after absorbing the bounded remainder.
\textbf{The exponential constant C is unchanged} --- $C=1/(2b_0)=24\pi^2/(11N_c)$ ---
but the displayed theorem should carry the \textbf{possibly reduced positive prefactor
c'\_0}. This is not an obstruction to the $e^{-C/g^2}$ behaviour; it corrects the
displayed statement, and the defect is in (18) alone. \emph{An earlier draft of this
note said the proof absorbed the remainder into the same symbol. It does not: that
reading came from a text extraction, in which the prime on $c_0'$ does not survive ---
a failure mode this programme has recorded before, and one to check against the
rendered page.}
\textbf{Theorem 6.3 retains the original H2 as a hypothesis}; a hypothesis that is not
discharged is still a hypothesis, and it may \emph{not} be presented as a theorem under
(H1) plus (H-FACT). But the statement is not literally established with its printed
quantifier over all $L$: the proof fixes $n_{\max}$ by the coupling and works on
lattices of extent $L/2^{n_{\max}}$, which is a lattice only for $L\in2^{n_{\max}}
\mathbb N$, or for other $L$ under an argument not given. \emph{An earlier draft of this
note called the statement untouched while conceding the missing condition; those cannot
both hold, and the second is the correct one.}

\emph{A citation loop is declared.} This paper's (H1) was said to be discharged by
\texttt{2602.0033}, whose Theorem 3.1 proves its terminal gap by ``applying the argument
of'' this paper's Theorem 5.3. Neither discharges a hypothesis of the other.

\section{Theorem 5.3: on a measure not identified with the physical one}

\subsection{Part (i): recoverable, but not by the printed proof}

Part (i) takes $H$ to be ``\emph{(the Friedrichs extension of) an elliptic operator with
smooth potential on $L^2(B)$}''. $L^2(B)$ is written without a measure; and $B=\A/\G$ is
a \emph{stratified space}, the reducible connections having stabilisers of positive
dimension, so ``the Friedrichs extension'' presupposes a closed form the paper does not
fix. This is the gap that forces the withdrawal of Corollaries 3.2--3.3 in
\texttt{ai.viXra:2602.0036}. \textbf{The repair does not need the quotient at all}: by
(I1) of the appended page, pullback is an isometry
$L^2(B,\nu)\cong L^2(\A,\mathrm{Vol}_{\A})^{\G}$ with $\nu=\pi_\#\mathrm{Vol}_{\A}$, and
$\A=SU(N_c)^E$ is a smooth compact connected manifold, where compact resolvent,
positivity improvement and uniqueness of the ground state are standard. It gives
$\Delta>0$ at fixed volume, qualitative only.

\subsection{Part (ii): the reference measure}

Step 1 puts $d\mu_0=\psi_0^2\,d\mathrm{Vol}_B$. The ground-state identity is correct and
is the right correction of v1, whose measure $e^{-2V_{\rm pot}/g^2}$ was dimensionally
inconsistent and failed numerically by more than fifty per cent. \textbf{But
$L^2(B,d\mathrm{Vol}_B)$ is not identified with the physical measure}: the physical
sector is $L^2(B,\nu)\cong L^2(\A)^{\G}$ by (I1). On the regular stratum
$d\nu=w\,d\mathrm{Vol}_B$, so the two measures are related by that density, and
\textbf{the associated generators differ by a drift}
$\langle\nabla\log w,\nabla\,\cdot\,\rangle$. They coincide only under an additional
constancy or identification statement for $w$, none of which is established here.
\emph{No inequality between the two is asserted}; what is asserted is that no
identification is available.

\emph{Two warnings about symbols, both load-bearing.} First, \textbf{the appended page
records this paper in two rows}: an \emph{as-printed} row carrying
$\mu_{0,\rm print}=\psi_{0,\rm print}^2\,d\mathrm{Vol}_B$, and a \emph{candidate
physical-restatement} row carrying $\mu_{0,\nu}=\psi_{0,\nu}^2\,d\nu$. An earlier
single-row version exposed the intended physical object but failed to record the
printed one, and is superseded. Second, \textbf{the $\nu$ of the appended page is not
the $\nu$ of the restored pages}: the page writes
$\nu_{\rm push}=\pi_\#\mathrm{Vol}_{\A}$, while Step 2 of Theorem 5.3 calls
$\nu_{\rm Vol}=d\mathrm{Vol}_B/\mathrm{Vol}_B(B)$ ``the uniform measure $\nu$''.
Identifying the two is precisely what (I2b) forbids.

\subsection{Part (ii): the base inequality}

Step 2 claims $\alpha(d\mathrm{Vol}_B)\ge N_c/2$ from $\operatorname{Ric}_B\ge N_c/4$ by
Bakry--\'Emery, hence $\lambda_1(d\mathrm{Vol}_B)\ge\alpha/2\ge N_c/4$. That pointwise
Ricci bound is Theorem 3.1 of \texttt{ai.viXra:2602.0036} and \textbf{is not
established}: its O'Neill trace omits mixed horizontal--vertical sectional terms which
for a bi-invariant metric are non-negative and subtract. \textbf{No pointwise Ricci
bound on the orbit space may presently be cited.}

\subsection{One repair for both --- and what it costs}

Work on $\mathcal H_{\rm phys}=L^2(\A,\mathrm{Vol}_{\A})^{\G}\cong L^2(B,\nu)$; take
$H_\nu=\frac{g^2}{2}K_\nu+V_{\rm pot}$ --- in flat notation, take H\_nu = (g\^2/2)
K\_nu + V\_pot, and \textbf{not} the gauge-fixed $V=S_{\rm YM}-\log\det M$, whose symbol
this manuscript has already fixed --- with $K_\nu$ the operator of the closed form on
$L^2(B,\nu)$; put $d\mu_0=\psi_{0,\nu}^2\,d\nu$. For the base inequality use
\texttt{ai.viXra:2602.0046} v3, which gives $\mathrm{RCD}^*(N_c/4,\dim\A)$ for
$(B,d_B,\nu)$, is valid across the singular locus, and \emph{explicitly bypasses} the
O'Neill computation. Holley--Stroock against $\nu$ then gives the rate.

\textbf{This is not free.} $H_\nu$ is a separately defined operator that
\textbf{has not been identified with the printed H}; correspondingly its ground state
$\psi_{0,\nu}$ has not been identified with the printed one. So the printed oscillation
hypothesis does not supply what the corrected route needs, and that route requires a
\emph{new} hypothesis, in flat notation (A-osc)\_nu:
\[
(\text{A-osc})_\nu:\qquad \operatorname{osc}\bigl(-\log\psi_{0,\nu}^2\bigr)\le
C\,L^d/g^2 .
\]
The route is worth recording --- it replaces two unavailable inputs with one available
one, and takes the quotient singularities out of the base inequality --- but it also
\emph{moves} the semiclassical hypothesis to a new object. It is not carried out here
and gives no uniformity in $L$. \textbf{The present fixed-volume argument does not by
itself establish (H1)}, and it has no bearing on (H2). \emph{An earlier draft of this
note said it bore in no way on (H1) either, while \S7 called it a route towards (H1);
\S7(B) now states what would have to be built for the two to meet.}

\section{What this paper does establish, stated exactly}

\begin{enumerate}\setlength{\itemsep}{0pt}
\item \textbf{The Hamiltonian/path-integral distinction} of \S7.2, in its narrow form:
the multiplication potential $V_{\rm pot}=S_{\rm YM}$ does not contain the gauge-fixed
term $-\log\det M$, which establishes a distinction between the two formulations. Its
third listed ingredient, $\operatorname{Ric}_B>0$, is not currently available, and this
narrow statement does not depend on it. \textbf{It does not establish that all
orbit-volume or Faddeev--Popov geometry is absent from the physical Hamiltonian}, whose
kinetic operator is $K_\nu$. \textbf{Orbit-volume geometry is present in the physical
measure and in the drift $\langle\nabla\log w,\nabla\,\cdot\,\rangle$ of $K_\nu$; the
Faddeev--Popov term is present in the gauge-fixed formulation; and this note does not
identify $w$ with $\det M$, so it does not claim that the same term has merely moved
from one formulation to the other.} \emph{Two earlier drafts erred in opposite
directions: the first certified the distinction unqualified, the second said the
geometry ``has moved from'' the potential --- itself an identification, asserted in the
sentence denying one.}
\textbf{And the measure formula printed alongside it is itself false: Remark 1.1
double-counts the Faddeev-Popov determinant.} It writes the gauge-fixed measure as
$e^{-V}\det M\,dA$ while setting $V=S_{\rm YM}-\log\det M$, which gives
$e^{-S_{\rm YM}}(\det M)^2dA$. The two correct forms, each with one determinant, are
\[
e^{-S_{\rm YM}}\det M\,dA
\qquad\text{or}\qquad
e^{-(S_{\rm YM}-\log\det M)}\,dA ,
\]
in flat notation exp(-S\_YM) det M dA and exp(-(S\_YM - log det M)) dA, and the printed
combination is neither. The narrow distinction between $V_{\rm pot}$ and
the gauge-fixed effective potential survives that correction.
Relatedly, Lemma 5.1 should not be read as an ordinary Hessian on $B$ ``at every point
of $\Mf$'': on the stratified quotient no such Hessian is fixed. What holds is that
\textbf{the lifted second variation on A is non-negative} --- every flat connection is a
global minimum of the Wilson action on the smooth $\A$ --- together with the
corresponding statement on $B_{\rm reg}$.
\item \textbf{Fixed-volume positivity}, once restated on $L^2(\A,\mathrm{Vol}_\A)^\G$.
Qualitative only.
\item \textbf{The exact ground-state-measure identity}, stated against the reference
measure it should carry. It is what gave this revision its reason to exist, and it
\textbf{invalidates the v1 derivation unless the required identification is proved};
\textbf{the one-plaquette suite supplies strong numerical counterevidence, not an
analytic refutation} --- the same rule this note applies to every other number the
suite produces. \emph{An earlier draft called it a refutation.}
\item \textbf{The constant} $C=24\pi^2/(11N_c)=1/(2b_0)$, resolving v1's internal
inconsistency between its introduction and its Theorem 6.3.
\item \textbf{Remark 5.4} as a description of the conditional bound: not uniform in $L$.
\item \textbf{Remark 5.6} in part: the one-plaquette evidence and the conclusion of
non-sharpness stand. Its inference that the gap on $B$ is $O(g)$ comes from Proposition
4.1, whose frequencies inherit the defects of \S2--\S3, and is \emph{not} certified here.
\item \textbf{The $SU(2)$, $2^3$, trivial-holonomy numerics} as evidence for that case.
\end{enumerate}

\section{Three lists, and they are not the same list}

\textbf{(A) The conditional implication of Theorem 6.3} needs: (H-BAL), that is
Balaban's Theorem 6.1, which its own statement assumes; a precise (H1), naming the
physical transfer operator, its fixed spatial slice and its Hilbert space; the original
(H2), or a valid derivation of it --- for which (H-FACT) is at most one input; and the
admissible-volume condition. \textbf{Theorem 6.3, with (H-BAL), (H1) and (H2) assumed
directly, uses nothing from list (C)}.

\textbf{(B) A candidate route towards (H1)}, and only that, would have to
\textbf{construct from S\_nmax a transfer operator on a fixed spatial slice}, specify
its Hilbert space and measure, identify its generator or its gap with the corrected
physical Hamiltonian $H_\nu$ of \S5.4, and control the effect of the $\epsilon_X$ terms
and their temporal nonlocality --- and only then would $(\text{A-osc})_\nu$ and
terminal-scale uniformity become relevant. \textbf{(A-osc)\_nu is specific to the
candidate route to (H1)}; it is not a hypothesis of Theorem 6.3, which assumes (H1)
outright. \textbf{The present fixed-volume route does not discharge (H1)}: it supplies
none of those identifications, it yields a bound at fixed volume where (H1) is a
statement at the terminal scale, and $H_\nu$ is a Wilson Hamiltonian on a spatial
configuration space rather than the transfer operator of a four-dimensional effective
action. \emph{An earlier draft listed $(\text{A-osc})_\nu$ among the requirements of
(A), and described (B) as though the two objects were already the same one.}

\textbf{(C1) The printed formulas that must be corrected}: Proposition 2.2(b),
equations (5) and (6), Proposition 3.3 and its Fourier labels, and equations (8)--(10).
This list is independent of (A). \emph{An earlier draft put (A) and (C) into one
sentence, which overstated what the conditional chain depends on.}

\textbf{(C2) What a Born--Oppenheimer theorem needs beyond those corrections}: a region
uniformly separated from the singular strata; a uniform normal gap \emph{on that
region}, which \S3.5 shows does not exist on the smooth stratum as a whole; control of
the reduction errors; and a construction that splices in or absorbs the quartic sector
on approach to the singular locus. \textbf{(C1) is necessary but not sufficient for a
rigorous Born-Oppenheimer theorem}, and the restored text concedes as much in asking
for a normal-gap condition it does not verify.

\section{Two blocks of the restored text are superseded wholesale}

Rather than chase each sentence: \textbf{Table 1 of Section 7.1 is superseded} by
\S\S3--7 of this note, and \textbf{the Note added in v2 is superseded} likewise.
Neither may be read as a current status summary. Table 1 still lists as ``gap proved''
the Paper II row resting on $\operatorname{Ric}_B\ge N_c/4$ and the Paper III row
resting on ground-state measure plus Holley--Stroock, both of which \S5 withdraws. The
Note added still carries item (1) on Morse--Bott and Proposition 3.3, item (4) on (H1)
discharged and (H2) refined, and \textbf{item (5), conclusion unaffected, is
withdrawn} together with them.

\emph{What survives from those two blocks}: the constant $C=24\pi^2/(11N_c)$ and the
resolution of v1's internal inconsistency about it; the repairs of broken citations,
of the non-existent ``Theorem 5.5'', and of the displaced headers; and \textbf{the
record that the v1 derivation is invalid without the missing identification, and that
the suite gives strong numerical counterevidence to its literal measure}
$e^{-2V_{\rm pot}/g^2}$. Everything else in Table 1 and in the Note added is to be
read through \S\S3--7 above.

\section{Scope}

Not audited here: the Balaban package (H-BAL); \textbf{the remainder of Proposition 4.1
beyond the two claims refuted in \S3.5}; and the numerical suite beyond the points
named. \emph{An earlier draft listed Proposition 4.1 as unaudited outright, which
contradicted \S3.5 of the same note} --- a partial audit of a proposition is not a
certification of the rest of it, and it is not a refusal to look either.

\vspace{3pt}
\noindent\rule{\textwidth}{0.8pt}

\noindent\footnotesize
\textbf{Provenance.} \S4 follows from the errata attached to \texttt{2602.0036} and
\texttt{2602.0033}. \S5.2 was found by asking which reference measure each paper writes
its ground-state measure against --- the question the appended page exists to force ---
and that page's single row for this paper exposed the intended physical object while
failing to record the printed one --- itself a provenance error, now split in two.
\textbf{\S2 came from an external audit of the second draft of this note}, which had
found three independent defects and still missed the angle ambiguity creating several
more; the first draft had opened by claiming nothing was false at all. \emph{Each round
has found the previous round's survivors list too short} --- the pattern this campaign
keeps measuring.

\end{document}
